% ============================================================
%  2.2  DNN Workloads and Notation
% ============================================================
\section{DNN Workloads and Notation}
\label{sec:bg-workloads}

Having established what a DNN is in Section~\ref{sec:bg-ai}, we now turn to
formalizing the \emph{computational workload} that each layer imposes on
hardware.
This formalization is essential because the mapping problem---how to schedule
and tile a workload onto a spatial architecture---is defined entirely in
terms of the loop-nest structure derived from it.


% ------------------------------------------------------------
\subsection{A Layer as a Loop Nest}
\label{sec:loop-nest}
% ------------------------------------------------------------

Every DNN layer can be decomposed into one or more
\emph{multiply-accumulate} (MAC) operations applied to three tensors:
an \emph{input activation} tensor, a \emph{weight} (or filter) tensor,
and an \emph{output activation} tensor.
A bias term, when present, is absorbed into the output initialization and
does not affect the loop-nest structure.

Because the MAC is associative and commutative, the order in which the
individual multiplications and accumulations are carried out is
irrelevant to the final numerical result.
This observation is crucial: it means the computation can be expressed as
a set of perfectly nested loops iterating over the various tensor
dimensions, and one is free to tile, reorder, and parallelize these loops
without altering correctness.

\paragraph{Example: one-dimensional convolution.}
To make the discussion concrete, consider a 1D convolution that maps an
input with $C$ channels and spatial extent $H$ to an output with $M$
channels and spatial extent $P$, using filters of size $R$:
\begin{equation}
\label{eq:1dconv}
\texttt{Output}[m,\, p] \;=\; \sum_{c=0}^{C-1}\;\sum_{r=0}^{R-1}
  \texttt{Input}[c,\, p{+}r]\;\times\;\texttt{Filter}[m,\, c,\, r],
\end{equation}
where $m \in [0,\,M)$ and $p \in [0,\,P)$.
Each output element requires $C \times R$ MAC operations, and the total
work is $M \times P \times C \times R$ MACs.

This expression can be read directly as a four-deep loop nest:

\begin{center}
\begin{minipage}{0.55\textwidth}
\begin{verbatim}
  for m in [0, M):
    for p in [0, P):
      for c in [0, C):
        for r in [0, R):
          Out[m][p] += W[m][c][r] * In[c][p+r]
\end{verbatim}
\end{minipage}
\end{center}

\noindent
The four loop dimensions---$M$, $P$, $C$, $R$---constitute the
\emph{iteration space} of the computation.
The key property is that these loops can be freely reordered (all $4!$
permutations produce the same output) and each loop can be tiled
(split into an outer and an inner sub-loop) to create smaller sub-problems
that fit in limited on-chip memories.
This flexibility is the foundation of all mapping strategies discussed in
Section~\ref{sec:bg-mapping}.

\paragraph{Extension to 2D convolutions.}
The standard 2D convolution, as used in CNNs, adds a second spatial
dimension and a batch dimension, leading to a seven-deep loop nest:
\begin{equation}
\label{eq:2dconv}
\texttt{O}[n,\,m,\,p,\,q]
  = \sum_{c}\;\sum_{r}\;\sum_{s}
    \texttt{I}[n,\,c,\,p{+}r,\,q{+}s] \;\times\;
    \texttt{W}[m,\,c,\,r,\,s],
\end{equation}
with dimensions $N$ (batch), $M$ (output channels), $P$ and $Q$ (output
spatial), $C$ (input channels), and $R$, $S$ (filter spatial).
This is the most common single-layer workload targeted by DNN accelerators
and the one used in the experimental evaluation of this thesis.


% ------------------------------------------------------------
\subsection{Coupling: the Relationship Between Dimensions and Tensors}
\label{sec:coupling}
% ------------------------------------------------------------

Each loop dimension in the nest participates in the indexing of one or
more of the three operand tensors.
The precise binding between dimensions and tensors---which we call the
\emph{coupling}---fully determines what data each iteration reads and
writes, and therefore what reuse opportunities exist.

\begin{definition}[Coupled and orthogonal dimensions]
\label{def:coupled-orthogonal}
A dimension $d$ is said to be \textbf{coupled} to an operand if $d$
appears in that operand's index expression (i.e., changing $d$ changes
the element accessed).
Conversely, $d$ is \textbf{orthogonal} to an operand if $d$ does not
appear in any of its index expressions; the operand's tile remains
constant as $d$ iterates, enabling complete \emph{stationarity} of that
operand.
\end{definition}

Consider the 2D convolution of Eq.~\eqref{eq:2dconv}.
Table~\ref{tab:coupling-conv} shows the coupling for each operand.
Every dimension is coupled to exactly two of the three operands and
orthogonal to one, creating a clean three-way partition that directly
maps onto the classical stationary dataflows (weight-stationary,
input-stationary, output-stationary), as discussed in
Section~\ref{sec:bg-mapping}.

\begin{table}[htbp]
\centering
\caption{Coupling of the standard 2D convolution.
A \checkmark\ indicates the dimension is coupled to (indexes) the operand;
a -- indicates orthogonality.}
\label{tab:coupling-conv}
\smallskip
\begin{tabular}{lccc}
\toprule
\textbf{Dimension} & \textbf{Input} & \textbf{Weight} & \textbf{Output} \\
\midrule
$M$ (out channels)   & --             & \checkmark      & \checkmark \\
$C$ (in channels)    & \checkmark     & \checkmark      & -- \\
$R$ (filter height)  & \checkmark     & \checkmark      & -- \\
$S$ (filter width)   & \checkmark     & \checkmark      & -- \\
$P$ (out height)     & \checkmark     & --              & \checkmark \\
$Q$ (out width)      & \checkmark     & --              & \checkmark \\
$N$ (batch)          & \checkmark     & --              & \checkmark \\
\bottomrule
\end{tabular}
\end{table}


% ------------------------------------------------------------
\subsection{Normal and Affine Dimensions}
\label{sec:normal-affine}
% ------------------------------------------------------------

Not all loop dimensions interact with tensor indices in the same way.
We distinguish two fundamental categories:

\begin{definition}[Normal dimension]
\label{def:normal-dim}
A dimension $d$ is \textbf{normal} if it appears alone as a tensor index
for every operand it is coupled to.
A normal dimension indexes each coupled operand independently: changing
$d$ selects a completely \emph{different} element along that operand's
corresponding axis.
\end{definition}

\begin{definition}[Affine dimension]
\label{def:affine-dim}
A dimension $d$ is \textbf{affine} if, for at least one operand, it
participates in a \emph{sum of indices} used to access a tensor
dimension.
That is, the operand is indexed by an expression of the form
$\alpha\, d + \beta \,d'$ (more generally, an affine combination of
two or more loop variables), rather than by $d$ alone.
\end{definition}

\paragraph{Illustration.}
In the 2D convolution of Eq.~\eqref{eq:2dconv}, the input
is indexed along its height by $p + r$.
Therefore:
\begin{itemize}
  \item $M$ and $C$ are \textbf{normal} dimensions: $M$ alone
    indexes the weight and output channel axes, $C$ alone indexes the
    input and weight channel axes.
  \item $R$ is \textbf{affine}: it appears in the sum $p + r$ that
    indexes the input's height.
    Likewise $S$ is affine via $q + s$.
  \item $P$ is also \textbf{affine}: it appears in the same sum $p + r$.
    $Q$ is affine via $q + s$.
\end{itemize}

The normal/affine distinction has deep consequences for mapping:
\begin{enumerate}
  \item \textbf{Tile-size computation.}
    When a normal dimension $d$ is tiled to size $t_d$ at some level, the
    operand tile along that axis is exactly $t_d$ elements.
    For an affine pair such as $(P,\,R)$, the input tile along the
    corresponding spatial axis is $t_P + t_R - 1$ (assuming unit stride),
    because consecutive output tiles \emph{overlap} in their input
    requirements.
  \item \textbf{Halo reuse.}
    When the innermost iterated dimension at a memory level is affine for
    an operand, consecutive iterations produce overlapping input tiles.
    Only the non-overlapping slice---the \emph{halo}---needs to be
    refetched, while the remaining data is reused from the previous
    iteration.
    This \emph{halo reuse} is unique to affine dimensions and does not
    arise with normal ones.
  \item \textbf{Map-space counting.}
    Only normal dimensions contribute independent tiling factors to the
    map-space cardinality (see Section~\ref{sec:bg-mapspace}).
    Affine dimensions affect tile sizes but are not independently looped;
    their tile sizes are derived from those of the normal dimensions
    they are summed with.
\end{enumerate}


% ------------------------------------------------------------
\subsection{Extended Einsum Notation}
\label{sec:een}
% ------------------------------------------------------------

To express the structure of tensor computations more concisely,
this thesis adopts the \textbf{extended Einsum notation}
(EEN)~\cite{abdelfattah2016EEN}, which augments the standard Einstein
summation convention with three features particularly suited to DNN
workloads:
\begin{enumerate}
  \item \textbf{Named ranks.}
    Each tensor dimension is given an explicit name (uppercase letter),
    and the variable iterating over it is denoted by the corresponding
    lowercase letter.
  \item \textbf{Shape tags.}
    Superscripts on each tensor specify the size of every rank, making
    the shapes explicit in the notation.
  \item \textbf{Affine subscripts.}
    Index expressions may contain sums of variables (e.g., $p + r$),
    directly encoding sliding-window and strided access patterns.
\end{enumerate}

Using the extended Einsum notation, the 1D convolution from
Eq.~\eqref{eq:1dconv} can be rewritten as:
\begin{equation}
\label{eq:een-1dconv}
\mathrm{Output}^{M,\,P}_{m,\,p}
  \;=\; \sum_{c,\,r}\;
  \mathrm{Input}^{C,\,H}_{c,\,p+r} \;\times\;
  \mathrm{Filter}^{M,\,C,\,R}_{m,\,c,\,r}\,.
\end{equation}

\noindent
The superscripts define the shapes of the tensors (e.g., the input has
$C$ channels and spatial extent $H = P + R - 1$).
The subscript $p + r$ makes the sliding-window access pattern
immediately visible: as $p$ advances, the accessed region in the input
shifts by one, overlapping with the previous region in $R - 1$ elements.
Padding is handled implicitly: any access with indices outside the
valid bounds of a rank is treated as zero.

The 2D convolution of Eq.~\eqref{eq:2dconv} becomes:
\begin{equation}
\label{eq:een-2dconv}
\mathrm{O}^{N,M,P,Q}_{n,m,p,q}
  \;=\; \sum_{c,r,s}\;
  \mathrm{I}^{N,C,H,W}_{n,\,c,\,p+r,\,q+s} \;\times\;
  \mathrm{W}^{M,C,R,S}_{m,c,r,s}\,.
\end{equation}

\noindent
The extended Einsum notation is a compact specification that maps
directly to the loop-nest representation of
Section~\ref{sec:loop-nest}.
Every rank in the notation corresponds to a loop dimension, every affine
subscript encodes a coupled (non-orthogonal) relationship between two
dimensions, and the implicit summation convention ($\sum_{c,r,s}$)
identifies the contraction (reduction) dimensions.
This one-to-one correspondence between the algebraic notation and the
loop-nest representation makes EEN a natural interface for specifying
workloads to mapping frameworks.


% ------------------------------------------------------------
\subsection{Activation-Dominant and Weight-Dominant Workloads}
\label{sec:act-weight-dominant}
% ------------------------------------------------------------

The coupling structure and dimension sizes of a layer determine which
operand generates the most data traffic.
This distinction is important for selecting mapping strategies and
for understanding when layer fusion is most beneficial.

\begin{definition}[Activation-dominant layer]
A layer is \textbf{activation-dominant} when the total size of its
activation tensors (input and output feature maps) exceeds the total
size of its weight tensor.
This is typical of the early layers of a CNN, where spatial dimensions
($P$, $Q$) are large and the number of channels ($M$, $C$) is small.
\end{definition}

\begin{definition}[Weight-dominant layer]
A layer is \textbf{weight-dominant} when its weight tensor is
larger than its activation tensors.
This occurs in later layers of deep networks, where spatial dimensions
have been reduced (through pooling or strided convolution) and the
channel count has grown substantially.
\end{definition}

\paragraph{Implications for layer fusion.}
Layer fusion primarily targets the elimination of intermediate
activation traffic between consecutive layers
(see Section~\ref{sec:bg-layer-fusion}).
Its benefit is therefore most pronounced in activation-dominant layers,
where the suppressed DRAM transfers of feature maps constitute a large
fraction of total data movement.
In weight-dominant layers, the potential savings from fusion are smaller
in relative terms, since the dominant traffic component---the
weights---is not reduced by fusion.
This trade-off is examined quantitatively in the experimental
evaluation (Chapter~\ref{ch:experimental-results}).
